% !TeX spellcheck = it_IT

\documentclass[format=169, handout]{beamer}

\usepackage[utf8]{inputenc}
\usepackage[italian]{babel}
\usepackage{multicol}

\usetheme{metropolis}

\usepackage{mathtools}
\DeclarePairedDelimiter{\ceil}{\lceil}{\rceil}

\title{Esercitazioni di Informatica B}
\subtitle{Array}

\author{Stefano Cereda\\
	stefano.cereda@polimi.it
}
\date{16/2018}
\institute[PoliMi]{\vspace{0.5cm}\centering Politecnico di Milano \\ \vspace{0.2cm}
	\includegraphics[width=0.2\textwidth]{../logopolimi}}

\setbeamercovered{invisible}

\makeindex

\begin{document}
	\begin{frame}
	\maketitle
	\end{frame}


\begin{frame}{Copia Array}
Scrivere un programma che, ricevuto un array $A$ di interi (con lunghezza nota $N$), copi il valore assoluto di ogni elemento di $A$ in un nuovo array $B$. 
\end{frame}

\begin{frame}
\frametitle{Conversione in maiuscolo}
Scrivere un programma che legga una sequenza di caratteri terminata dal carattere '0' e lunga al massimo N caratteri.
Il programma trasforma tutte le lettere alfabetiche minuscole nelle corrispondenti lettere maiuscole.

Si ricordi che la codifica ASCII mette in corrispondenza i caratteri con dei numeri.

In seguito, si modifichi il programma in modo da sostituire tutti i caratteri punto, virgola, due punti, punto e virgola e punto di domanda in spazi.
\end{frame}

\begin{frame}{Ricerca prima occorrenza - esempio da lezione}
Scrivere un programma che legga un vettore di interi $V$ ed un numero intero $n$ per poi stampare l'indice di $V$ corrispondente al primo elemento uguale ad $n$.
\end{frame}

\begin{frame}{Ricerca prima occorrenza - estensioni}
\footnotesize
Dato un vettore ordinato di interi $V$ ed un numero $n$, si scriva la porzione di codice per implementare un algoritmo di ricerca e stampare l'indice del primo elemento di $V$ uguale ad $n$. (Esercizio precedente)

Si modifichi il programma per trovare più occorrenze dello stesso numero (stampandone tutti gli indici).

\pause
Si modifichi il programma per creare un secondo array $V_2$ in cui inserire gli indici delle occorrenze (quelli stampati al punto precedente). Si presti attenzione a non lasciare buchi in $V_2$. Al termine si stampi $V_2$.

\pause
Ad esempio:
\[ V_1 = [1, 3, 5, 1, 4, 6, 1, 8, 1] \]
\[ n=1 \]
\[ V_2 = [0, 3, 6, 8]\]
E non:
\[ V_2 = [0, ?, ?, 3, ?, ?, 6, ?, 8]\]
\end{frame}

\begin{frame}{Conversione binaria}
Si scriva un programma che consenta di trasformare un numero da base 10 a base 2 usando l'algoritmo delle divisioni ripetute. Si stampi il risultato nell'ordine corretto. Si effettui la conversione con un numero massimo di bit $N$.
\end{frame}

\begin{frame}{Radice}
Scrivere un programma che riceva dall'utente 20 numeri float positivi.
Si controlli che l'utente non inserisca numeri negativi, nel caso si chiede di effettuare nuovamente l'inserimento.
Il programma memorizza in un array le radici quadrate dei numeri inseriti.
\end{frame}

\begin{frame}{Ricerca per bisezione}
Si scriva un programma che riceva in ingresso un vettore ordinato ed un valore.
Il programma deve cercare il valore all'interno dell'array usando l'algoritmo di ricerca per bisezione.
\end{frame}

\begin{frame}{Ordinamento a bolle}
Si scriva un programma che riceva in ingresso un vettore di numeri interi.
Il programma deve ordinare il vettore usando l'algoritmo dell'ordinamento a bolle e stampare il risultato.
\end{frame}
\end{document}